\chapter{Memoria de calculo de iluminacion}
\label{chap:iluminacion}

\section{Objetivo y metodologia}
Esta memoria determina la cantidad de luminarias, la potencia instalada y la
densidad de potencia de iluminacion (LPD) de cada ambiente del grifo mediante
el \textbf{metodo de lumanes (metodo de rendimiento o de factor de
utilizacion)}. El procedimiento es el siguiente:

\begin{enumerate}
\item Definir el nivel de iluminancia objetivo \(E\) segun el uso del ambiente,
  con base en la NTP ISO 8995 y la EM.010 del RNE.
\item Calcular el \textbf{indice de local} \(K\), que relaciona las
  dimensiones del ambiente con la altura de montaje sobre el plano de trabajo:
  \[
    K = \frac{L \cdot A}{H_m\,(L + A)}
  \]
  siendo \(L\) el largo, \(A\) el ancho y \(H_m\) la altura de montaje
  (altura del local menos la altura del plano de trabajo).
\item Adoptar el \textbf{factor de utilizacion} \(F_U\), estimado por
  interpolacion de tablas tipicas de luminarias LED segun \(K\) y las
  reflectancias del ambiente.
\item Adoptar el \textbf{factor de mantenimiento} \(F_M\) (\(0,80\) en
  interiores limpios y \(0,70\) en exteriores), que considera el ensuciamiento
  de la luminaria y la depreciacion del flujo luminoso.
\item Calcular el \textbf{flujo luminoso requerido}:
  \[
    \Phi = \frac{E \cdot S}{F_U \cdot F_M}
  \]
  siendo \(S\) el area del ambiente.
\item Determinar la cantidad de luminarias:
  \[
    N = \left\lceil \frac{\Phi}{\phi_l} \right\rceil
  \]
  con \(\phi_l\) el flujo luminoso de la luminaria seleccionada.
\item Calcular la \textbf{potencia instalada}, la \textbf{densidad de potencia
  de iluminacion} \(\text{LPD} = P/S\) y la \textbf{iluminancia resultante}:
  \[
    E_{res} = \frac{N \cdot \phi_l \cdot F_U \cdot F_M}{S} \geq E
  \]
\end{enumerate}

Todos los valores son un criterio academico de anteproyecto. La seleccion
final de luminaria se confirma con catalogo vigente y la verificacion de campo.

\section{Parametros de diseno}
\begin{itemize}
\item Niveles de iluminancia segun NTP ISO 8995 y EM.010 (oficinas 300--500~lux,
  sala de maquinas 200~lux, servicios higienicos 150~lux, zona de despacho
  150--200~lux, circulacion exterior 50~lux).
\item Factor de mantenimiento \(F_M=0,80\) interior y \(F_M=0,70\) exterior.
\item Luminarias LED de 100~lm/W (panel de oficina 36~W, panel SS.HH 18~W,
  high-bay industrial 50~W, proyector exterior 100~W, poste de patio 80~W),
  referenciales de familias de catalogo comparables.
\item Reflectancias tipicas: techo 0,7, pared 0,5 y piso 0,2 en interiores;
  reflectancia exterior 0,1.
\end{itemize}

\section{Ejemplo de calculo completo}
Se desarrolla el ambiente \textbf{\IllumEjNombre} con las siguientes entradas:

\begin{center}
\small
\begin{tabularx}{\linewidth}{@{}l l X@{}}
\toprule
Parametro & Valor & Fuente / nota \\
\midrule
Dimensiones & \IllumEjLargo~m \(\times\) \IllumEjAncho~m \(\times\) \IllumEjAltura~m &
  arquitectura extraida del DWG (criterio adoptado) \\
Area & \IllumEjArea~m\textsuperscript{2} & \(S = L \cdot A\) \\
Iluminancia objetivo & \IllumEjLux~lux & NTP ISO 8995 \\
Altura de montaje & \IllumEjAltura~m \(-\,0,85\)~m = 2,15~m & plano de trabajo 0,85~m \\
\bottomrule
\end{tabularx}
\end{center}

\textbf{Indice de local.}
\[
K = \frac{L \cdot A}{H_m\,(L+A)}
  = \frac{\IllumEjLargo \cdot \IllumEjAncho}{2,15\,(\IllumEjLargo + \IllumEjAncho)}
  = \IllumEjK
\]

\textbf{Factores.} Para \(K = \IllumEjK\) y reflectancias tipicas de oficina se
adopta \(F_U = \IllumEjFu\); el factor de mantenimiento es \(F_M = \IllumEjFm\).

\textbf{Flujo luminoso requerido.}
\[
\Phi = \frac{E \cdot S}{F_U \cdot F_M}
     = \frac{\IllumEjLux \cdot \IllumEjArea}{\IllumEjFu \cdot \IllumEjFm}
     = \IllumEjFlujoReq~\text{lm}
\]

\textbf{Cantidad de luminarias.} Con la luminaria seleccionada (panel LED de
\IllumEjLumFlujo~lm y \IllumEjLumPot~W):
\[
N = \left\lceil \frac{\IllumEjFlujoReq}{\IllumEjLumFlujo} \right\rceil
  = \IllumEjN~\text{luminarias}
\]

\textbf{Potencia instalada y densidad de potencia.}
\[
P = N \cdot P_l = \IllumEjN \cdot \IllumEjLumPot~\text{W} = \IllumEjPot~\text{W},
\qquad
\text{LPD} = \frac{P}{S} = \frac{\IllumEjPot}{\IllumEjArea} = \IllumEjLpd~\text{W/m}^2
\]

\textbf{Verificacion del nivel de iluminacion obtenido.}
\[
E_{res} = \frac{N \cdot \phi_l \cdot F_U \cdot F_M}{S}
        = \frac{\IllumEjN \cdot \IllumEjLumFlujo \cdot \IllumEjFu \cdot \IllumEjFm}{\IllumEjArea}
        = \IllumEjResult~\text{lux} \geq \IllumEjLux~\text{lux}
        \quad \checkmark
\]

El nivel resultante supera el objetivo, por lo que el ambiente CUMPLE.

\section{Resultados por ambiente}
\IllumNotaN

\ILLUMTAB

\section{Verificacion de densidad de potencia (LPD)}
La potencia total de iluminacion calculada por el metodo de lumanes es
\SI{\IllumTotalKW}{kW} (\IllumTotalN~luminarias), con una densidad promedio de
\SI{\IllumTotalLPD}{W/m^2}. Todos los ambientes cumplen los limites practicos
de LPD para iluminacion LED y el objetivo de iluminancia se alcanza en todos
los casos.

Estas potencias son coherentes con los circuitos de alumbrado del cuadro de
cargas: L-01 = 800~W (marquesina y despacho), L-02 = 640~W (postes de patio),
A1-01 = 540~W (administracion y oficina) y A1-02 = 222~W (sala de maquinas,
SS.HH 1--3 y vereda/venteo). El circuito L-03 corresponde al letrero/totem de
precios (400~W) y se trata como carga de rotulo luminoso, no como parte del
calculo de iluminancia por ambiente.

\section{Criterios normativos}
\begin{itemize}
\item NTP ISO 8995:2008, Iluminacion de lugares de trabajo (equivalente a la
  CIE S 008), de la que se toman los niveles de iluminancia por uso.
\item EM.010 del RNE, articulo 6, niveles de iluminancia, y articulo 11.1,
  alumbrado de emergencia y senales de evacuacion (tratado en el capitulo de
  calculos y especificaciones).
\item Practica recomendada de estaciones de servicio para la zona de despacho
  (150--200~lux) y la circulacion exterior (50~lux).
\end{itemize}

